% !TeX program = xelatex
\documentclass[aspectratio=169]{beamer}

\usepackage{import}

\subimport{../../}{system.tex}
\UseTemplateSet{
  layout = beamer,
  globalfonts = {cmu,shanggu},
  features = {math}
}

\title{世界文字總覽~導論}
\author{Kelvin Yang}
\date{\today}

\begin{document}
\frame{\titlepage}

\begin{frame}{說明}
    本課件只爲講解方便，\textbf{不}涵蓋講義的全部內容且往往\textbf{預設}某些知識，這些知識或可以從講義中得到補充。
\end{frame}

\begin{frame}{書寫系統與文字}
    區分以下兩個概念是有必要的：\textbf{書寫系統}（Writing System）與\textbf{文字}（script）。

    \begin{block}{書寫系統}
        能夠按照一定次序（往往是線性的）書寫某門語言的視覺符號。
    \end{block}

    \begin{block}{文字}
        相互之間十分「相似」的書寫系統。這種相似性往往意涵形式、使用邏輯、語音對應等，但更多是習慣上的。
    \end{block}

    對於語言$L\in\mathcal{L}$，其書寫系統$W\in\mathcal{W}$是一種結構，滿足結構上的某種類同構：
    \[
        L\simeq W.
    \]
    定義書寫系統的相似爲等價關係$\sim$，則文字定義爲等價類$\mathcal{S}:=\mathcal{W}/\sim$.
\end{frame}

\begin{frame}{文字類型學}
    按照字位與語言之間的映射關係，我們言及\textbf{文字類型學}（typology of scripts）。

    須注意：我們對文字類型的分析往往存在某種\textbf{簡化}，即有意忽略了其中相當一部分。這種簡化是必要的，否則會使分析過於複雜。

    \begin{block}{文字類型}
        \begin{itemize}
            \item \textbf{表音文字}（phonogram）：字位映射到態位以下
            \begin{itemize}
                \item \textbf{音段文字}（segmental）：字位映射到音段成分
                \begin{itemize}
                    \item \textbf{全音位文字}（alphabet）：字位一一映射到音位
                    \item \textbf{輔音音位文字}（abjad）：字位映射到音位，但輔音地位高於元音
                    \item \textbf{元音附標文字}（abugida）：字位映射到音位，但某個元音低位低於其他元音
                \end{itemize}
                \item \textbf{音節文字}（syllabary）：字位映射到音節
            \end{itemize}
            \item \textbf{表意文字}（ideogram）：\textbf{存在}能夠表意的字位
            \begin{itemize}
                \item \textbf{語素文字}（logogram）：\textbf{存在}字位映射到態位及以上
            \end{itemize}
        \end{itemize}
    \end{block}
\end{frame}

\begin{frame}{文字發展類型}
    我們將文字按其發展簡單分爲兩類：\textbf{獨立型}（independent type）、\textbf{他源型}（externally derived type）。

    \begin{block}{獨立型}
        自我發源後幾乎不受外來文字影響，相對獨立發展。也即\textbf{自源型}（primary type）。
    \end{block}

    \begin{block}{他源型}
        來源於既有的其他文字，可以再分爲：
        \begin{itemize}
            \item \textbf{適應型}（adapted type）：未產生新文字，總體上是把現有文字拿來用，只有少量創新
            \item \textbf{衍生型}（derived type）：在其他文字的啓發或影響下，創製出的新文字
        \end{itemize}
    \end{block}
    
    另外，我們也歸納出\textbf{競爭型}（competitive type），用以描述同一種語言（族羣內）存在多種文字並存競爭的情況。但這種情況要麼融合出新文字（衍生型），要麼共存（兩套獨立文字），要麼混合（類似適應型）。
\end{frame}

\begin{frame}{四大文字系統}
    我們嘗試將世界主要文字分爲四大系統：\textbf{希臘系}（Hellenic）、\textbf{漢字系}（Sinographic）、\textbf{阿拉米系}（Aramaic）和\textbf{婆羅米系}（Brahmic）。

    這個分系主要考慮文化影響力與文字發展進程，但實際上並不窮盡所有文字，且希臘、阿拉米、婆羅米三系很可能存在同源關係，漢字系也不是與這三者毫無來往的——總之，這只是一個簡化的框架。

    \begin{block}{希臘系}
        希臘系是目前影響力最爲廣泛的文字系統。希臘人藉由腓尼基文創造出了希臘文，將原本的輔音音位文字改造爲全音位文字，影響深遠。
        西支傳入義大利，並最後進入羅馬，最後產生了拉丁文。東支則傳到俄羅斯，成爲了現代西里爾文的源頭。
    \end{block}

    \begin{block}{漢字系}
        漢字系在東亞的影響最爲廣泛。漢字本身屬於相當典型的獨立型文字，在內啓發了諸多少數民族文字的創制，在外則傳播到日本、朝鮮、越南等諸多地區。
    \end{block}
\end{frame}

\begin{frame}{四大文字系統}
    \begin{block}{阿拉米系}
        阿拉米文最初是西亞地區阿拉米人所使用的文字，應當也來源於腓尼基文，在阿契美尼德王朝時期得到推廣。
        其發展大概可以分成三系：
        \begin{itemize}
            \item 西亞子系：衍生出希伯來文和敘利亞文
            \item 阿拉伯子系：衍生出納巴泰文，最後改造爲阿拉伯文，並隨着伊斯蘭教被傳播到廣大地區
            \item 中亞子系：敘利亞文東傳後衍生出粟特文、回鶻文，並發展爲蒙古文、滿文等
        \end{itemize}
    \end{block}

    \begin{block}{婆羅米系}
        婆羅米文是印度最古老的文字之一，屬於元音附標文字，發展出了現在南亞、東南亞等地區的絕大多數文字。
        北支由笈多文演變出天城體、藏文、孟加拉文、古吉拉特文等，南支則衍生出泰米爾文、泰盧固文、卡納達文等南印度文字以及東南亞的緬甸文、高棉文等諸多文字。
    \end{block}
\end{frame}

\begin{frame}{總綱}
    我們的總思路在於：以我們上述的四大系文字爲主體，分別介紹其歷史，並選取其中最重要的幾種文字來進行介紹與學習。

    對文字的學習，我們會選取一種最經典、規整的語言，僅僅是爲了可以讀出這些文字，以及方便記憶。
    但是，我們的介紹\textbf{不會}重視語言，這意味着我們會忽略語法、句法、詞彙等一切要素，而\textbf{僅僅}着眼在認讀與拼讀文字上面。
\end{frame}

\end{document}
